% =============================================================================
% Chapter 9 --- Conclusion
% Adapted from Context 5 §7 closing statement.
% =============================================================================

\chapter{Conclusion}
\label{ch:conclusion}

This report has presented the RTL design and functional verification of
a controller for an ANN accelerator with in-memory computing
peripherals, as implemented in the repository under the Technion VLSI
Laboratory course framework. The design comprises three synthesizable
SystemVerilog modules --- \modname{parallel\_interface},
\modname{ann\_controller}, and \modname{input\_buffer} --- together
with a shared package layer that centralizes the address-packing,
one-hot predicate, and pulse-train arithmetic conventions. The
controller is built around a hierarchical finite-state machine with
three cooperating sub-FSMs for programming, verification, and erase,
and around a data-driven LUT-based pulse shaper that tunes the first
programming attempt per expected weight.

The controller and its verification together establish, under explicit
modeling assumptions, a coherent digital control layer: command
decoding, one-hot addressing, burst-shaped pulses, program /
verify / erase recovery when the readback is \emph{forced} to misbehave,
and interoperability with the parallel interface and the input
buffer. The directed regression under \sv{target/} records 30 PASS
cases for the static address/pulse contract, 10 PASS of 10 for
\sv{PI}\,\(+\)\,controller\,\(+\)\,buffer integration, 8 PASS of 8 for
reorder-under-reads, successful host-erase selectivity against the
behavioral mock, and zero end-of-run mismatches across 640 programmed
weights, the last of which also exercises both the re-PROG and ERASE
recovery branches explicitly (14 ERASE triggers and 3017 re-PROG
events under the fault-injection strategy described in
\Cref{ch:verification}).

The same evidence makes the limits of the work explicit: behavioral
mocks stand in for the physical memristor array; \signame{op\_done} and
\signame{weight\_read\_data} are idealized; several compatibility hooks
in the RTL (\stateval{S\_RESET}, unused module parameters, the internal
\signame{weight\_read\_en}) are not part of the minimal proven subset;
and no netlist, timing closure, or power/area analysis is in scope. A
rigorous next phase is therefore not merely ``more tests'', but
assertions on the real protocols, randomized legal stimuli with
coverage of FSM and recovery branches, and --- where the program
requires it --- models that lift the abstraction from 4-bit tables to
the actual device non-idealities that the silicon team must sign off
on~\cite{Ielmini2018RRAM, Sebastian2020IMC, Wang2022YFlash}.

That progression stays faithful to the boundary of what the repository
can claim today while giving a VLSI follow-on team a defensible map
from the present RTL to a production-grade sign-off process.
